\documentclass[11pt,aspectratio=169]{beamer}

% 直接复用现有主题配置。
% 通用宏包集中在此处；课程内容文件中通常不应再加载宏包。
\usepackage{iftex}
\ifPDFTeX
  \usepackage[T1]{fontenc}
  \usepackage[utf8]{inputenc}
  \usepackage{lmodern}
\else
  \usepackage{fontspec}
  \defaultfontfeatures{Ligatures=TeX}
\fi

\usepackage{amsmath,amssymb,mathtools}
\usepackage{booktabs}
\usepackage{graphicx}
\usepackage{tikz}
\usetikzlibrary{arrows.meta,positioning,shapes.geometric}
\usepackage[most]{tcolorbox}
\usepackage{listings}
\usepackage{xparse}
\usepackage{appendixnumberbeamer}
\usepackage[backend=biber,style=authoryear,maxcitenames=2]{biblatex}
% 修改整套课件配色时，只需要编辑本文件。
\definecolor{Primary}{HTML}{003B6F}
\definecolor{PrimaryLight}{HTML}{EAF2F8}
\definecolor{Accent}{HTML}{D97706}
\definecolor{AccentLight}{HTML}{FFF4E5}
\definecolor{Background}{HTML}{FFFFFF}
\definecolor{Text}{HTML}{20252B}
\definecolor{MutedText}{HTML}{66717E}
\definecolor{RuleColor}{HTML}{D7DEE5}
\definecolor{Success}{HTML}{237A57}
\definecolor{SuccessLight}{HTML}{EAF6F0}
\definecolor{CodeBackground}{HTML}{F5F7F9}
% 字体语义集中定义，避免在每一页手动设置字号与字重。
\ifPDFTeX
  \usefonttheme{professionalfonts}
\else
  \IfFontExistsTF{TeX Gyre Heros}{
    \setsansfont{TeX Gyre Heros}
  }{
    \setsansfont{Latin Modern Sans}
  }
  % TeX Gyre Cursor 随常见 TeX Live 安装提供，适合作为可移植的代码字体。
  \setmonofont{TeX Gyre Cursor}
\fi

\newcommand{\TitleFont}{\fontsize{25}{30}\selectfont\bfseries}
\newcommand{\SubtitleFont}{\large}
\newcommand{\FrameTitleFont}{\large\bfseries}
\newcommand{\BodyFont}{\normalsize}
\newcommand{\SmallFont}{\small}
\newcommand{\CodeFont}{\ttfamily\small}
% Beamer 外观：白底、蓝色主色、无导航按钮、极简页脚。
\mode<presentation>

\useinnertheme{rectangles}
\usefonttheme{professionalfonts}
\setbeamertemplate{navigation symbols}{}
\setbeamertemplate{headline}{}
\setbeamersize{text margin left=0.75cm,text margin right=0.75cm}

\setbeamercolor{normal text}{fg=Text,bg=Background}
\setbeamercolor{structure}{fg=Primary}
\setbeamercolor{alerted text}{fg=Accent}
\setbeamercolor{frametitle}{fg=Primary,bg=Background}
\setbeamercolor{title}{fg=Primary}
\setbeamercolor{subtitle}{fg=MutedText}
\setbeamercolor{author}{fg=Text}
\setbeamercolor{institute}{fg=MutedText}
\setbeamercolor{date}{fg=MutedText}
\setbeamercolor{itemize item}{fg=Primary}
\setbeamercolor{itemize subitem}{fg=Accent}
\setbeamercolor{section in toc}{fg=Text}
\setbeamercolor{section in toc shaded}{fg=MutedText}

\setbeamerfont{frametitle}{series=\bfseries,size=\large}
\setbeamerfont{title}{series=\bfseries,size=\LARGE}
\setbeamerfont{footline}{size=\scriptsize}

\setbeamertemplate{frametitle}{%
  \vspace{0.18cm}%
  \insertframetitle\par
  \vspace{0.10cm}%
  {\color{RuleColor}\rule{\textwidth}{0.5pt}}%
}

\setbeamertemplate{footline}{%
  \leavevmode
  \hbox{%
    \begin{beamercolorbox}[wd=.78\paperwidth,ht=2.8ex,dp=1.2ex,leftskip=.75cm]{normal text}%
      \color{MutedText}\insertshorttitle\enspace |\enspace\insertshortauthor
    \end{beamercolorbox}%
    \begin{beamercolorbox}[wd=.22\paperwidth,ht=2.8ex,dp=1.2ex,rightskip=.75cm plus1fil]{normal text}%
      \hfill\color{MutedText}\insertframenumber\,/\,\inserttotalframenumber
    \end{beamercolorbox}%
  }
}

\setbeamertemplate{title page}{%
  \vfill
  {\color{Primary}\rule{1.2cm}{3pt}\par}
  \vspace{0.55cm}
  {\TitleFont\color{Primary}\inserttitle\par}
  \vspace{0.18cm}
  {\SubtitleFont\color{MutedText}\insertsubtitle\par}
  \vspace{0.9cm}
  {\BodyFont\insertauthor\par}
  {\SmallFont\color{MutedText}\insertinstitute\par}
  \vspace{0.35cm}
  {\SmallFont\color{MutedText}\insertdate\par}
  \vfill
}

\setbeamertemplate{section page}{%
  \vfill
  {\color{MutedText}\large Chapter \insertsectionnumber\par}
  \vspace{0.25cm}
  {\color{Primary}\rule{1.2cm}{3pt}\par}
  \vspace{0.45cm}
  {\fontsize{27}{32}\selectfont\bfseries\color{Primary}\insertsectionhead\par}
  \vfill
}

\AtBeginSection[]{
  \begin{frame}[plain,noframenumbering]
    \sectionpage
  \end{frame}
}

\mode<all>
% 页面级语义环境。正文优先使用这些环境，而不是直接使用 frame。
\newenvironment{titleframe}
  {\begin{frame}[plain,noframenumbering]\titlepage}
  {\end{frame}}

\newenvironment{outlineframe}[1][Outline]
  {\begin{frame}{#1}\tableofcontents[currentsection,hideothersubsections]}
  {\end{frame}}

\newenvironment{theoryframe}[1]
  {\begin{frame}{#1}}
  {\end{frame}}

\newenvironment{exampleframe}[1]
  {\begin{frame}{#1}}
  {\end{frame}}

\newenvironment{figureframe}[1]
  {\begin{frame}{#1}\centering}
  {\end{frame}}

% fragile 页面必须通过 environment=codeframe 告知 Beamer 外层环境名。
\newenvironment{codeframe}[1]
  {\begin{frame}[fragile,environment=codeframe]{#1}}
  {\end{frame}}

\newenvironment{summaryframe}[1][Summary]
  {\begin{frame}{#1}}
  {\end{frame}}

\newcommand{\figurecaption}[1]{%
  \par\vspace{0.35cm}{\small\color{MutedText}#1\par}%
}
% 内容组件：信息框、强调命令、代码样式与常用数学命令。
\tcbset{
  coursebox/.style={
    enhanced,
    boxrule=0pt,
    leftrule=2.5pt,
    arc=1pt,
    outer arc=1pt,
    left=7pt,right=7pt,top=6pt,bottom=6pt,
    fonttitle=\bfseries,
    before skip=6pt,after skip=6pt
  }
}

\newtcolorbox{definitionbox}[1][Definition]{
  coursebox,title=#1,colback=PrimaryLight,colframe=Primary,coltitle=Primary
}
\newtcolorbox{theorembox}[1][Theorem]{
  coursebox,title=#1,colback=PrimaryLight,colframe=Primary,coltitle=Primary
}
\newtcolorbox{proofbox}[1][Proof idea]{
  coursebox,title=#1,colback=SuccessLight,colframe=Success,coltitle=Success
}
\newtcolorbox{remarkbox}[1][Remark]{
  coursebox,title=#1,colback=CodeBackground,colframe=MutedText,coltitle=Text
}
\newtcolorbox{warningbox}[1][Warning]{
  coursebox,title=#1,colback=AccentLight,colframe=Accent,coltitle=Accent
}
\newtcolorbox{importantbox}[1][Important]{
  coursebox,title=#1,colback=AccentLight,colframe=Accent,coltitle=Accent
}

\newcommand{\keyterm}[1]{\textcolor{Primary}{\textbf{#1}}}
\newcommand{\coursecheck}{\textcolor{Success}{\ensuremath{\checkmark}}}

\lstdefinestyle{coursecode}{
  basicstyle=\CodeFont,
  backgroundcolor=\color{CodeBackground},
  keywordstyle=\color{Primary}\bfseries,
  commentstyle=\color{Success},
  stringstyle=\color{Accent},
  numbers=left,
  numberstyle=\tiny\color{MutedText},
  numbersep=8pt,
  frame=single,
  rulecolor=\color{RuleColor},
  framesep=7pt,
  showstringspaces=false,
  breaklines=true,
  columns=fullflexible,
  keepspaces=true,
  tabsize=2
}
\lstset{style=coursecode}

\newcommand{\R}{\mathbb{R}}
\newcommand{\N}{\mathbb{N}}

\addbibresource{iv-references.bib}
\DeclareMathOperator{\Cov}{Cov}
\DeclareMathOperator{\Var}{Var}
\newcommand{\E}{\mathbb{E}}
\newcommand{\Prob}{\mathbb{P}}
\tcbset{coursebox/.append style={colbacktitle=Background}}

% Automatic transition slide for each subsection:
% large section title above, subsection title below.
\newcommand{\sectionsubsectionframe}{%
  \begin{frame}[plain,noframenumbering]
    \vspace{0.9cm}
    {\fontsize{30}{36}\selectfont\bfseries\color{Text}%
      \insertsectionhead\par}

    \vspace{0.45cm}
    \noindent
    {\color{Accent}\rule{1.2cm}{1.3pt}}%
    {\color{RuleColor}\rule{\dimexpr\textwidth-1.2cm\relax}{0.6pt}}

    \vspace{0.45cm}
    {\fontsize{22}{27}\selectfont\bfseries\color{Text}%
      \insertsubsectionhead\par}
    \vfill
  \end{frame}%
}

\AtBeginSubsection[]{\sectionsubsectionframe}
% 课程信息：以后通常只需要修改这里和正文内容。
\title[Causal Inference]{Causal Inference}
\subtitle{Instrumental Variables}
\author[Zhaoyi]{Zhaoyi}
%\institute[Your University]{Your University}
\date{\today}

\begin{document}








\begin{titleframe}


\end{titleframe}


\begin{outlineframe}

\end{outlineframe}










\section{Randomized Controlled Trial}
\begin{frame}
A Randomized Controlled Trial (RCT) works as follows:
\begin{enumerate}
    \item Choose how to randomize (mechanism), e.g.:
        \begin{itemize}
            \item Bernoulli trial
            \item randomized experiment
            \item Stratified randomized experiment.
        \end{itemize}
    \item Pick $n$ units from the target population and observe pre-treatment characteristics $X_i$
    \item Randomly assign the treatment $\{D_i\}_{i=1}^n$ following the assignment mechanism
    \item Observe outcomes:$$\{Y_i\}_{i=1}^n \quad \text{with}\quad Y_i=D_iY_i(1)+(1-D_i)Y_i(0)$$
\end{enumerate}


\end{frame}



\section{Randomization Schemes}
\subsection{Bernoulli Trials}

\begin{frame}{Unbiasedness: Bernoulli Trial}
\begin{itemize}
    \item Consider a simple Bernoulli trial, where each unit is assigned to treatment independently of the others.
    \item Specifically, suppose treatment is randomized such that, conditional on all potential outcomes $\{Y_i(0), Y_i(1)\}_{i=1}^n$ and the number of treated units $n_1$, each unit has the same probability of receiving treatment:
    $$\mathbb{E}[D_i \mid \{Y_i(0), Y_i(1)\}_{i=1}^n, n_1] = \frac{n_1}{n}.$$
\end{itemize}

\end{frame}


\begin{frame}{Theorem: Unbiasedness under Random Assignment}
\begin{theorembox}[Unbiasedness under Randomized Assignment]
Under this framework,
$$\mathbb{E}\left[\widehat{\tau}_{DM} \mid \{Y_i(0), Y_i(1)\}_{i=1}^n, n_0 > 0, n_1 > 0\right] = \tau_{SATE}.$$
\end{theorembox}
That is, $\widehat{\tau}_{DM}$ is an unbiased estimator of the Sample Average Treatment Effect (SATE), even in finite samples.
\end{frame}

\begin{frame}{Theorem: Asymptotic Distribution of Difference in Means Estimator}

\begin{theorembox}[Theorem: Asymptotic Distribution of Difference in Means Estimator]
Suppose we are in the framework discussed above. Then,
$$\sqrt{n}(\widehat{\tau}_{DM} - \tau_{ATE}) \xrightarrow{d} \mathcal{N}(0, V_{DM}), \quad \text{where} \quad V_{DM} = \frac{\text{Var}[Y_i(0)]}{1 - \pi} + \frac{\text{Var}[Y_i(1)]}{\pi}.$$

Furthermore, the plug-in variance estimator
$$\widehat{V}_{DM} := \frac{n}{n_0^2} \sum_{i:D_i=0} \left(Y_i - \frac{1}{n_0} \sum_{j:D_j=0} Y_j\right)^2 + \frac{n}{n_1^2} \sum_{i:D_i=1} \left(Y_i - \frac{1}{n_1} \sum_{j:D_j=1} Y_j\right)^2$$
is consistent: $\widehat{V}_{DM} \xrightarrow{p} V_{DM}$.
\end{theorembox}
\end{frame}




\begin{frame}{OLS estimator and RCTs}

\begin{itemize}
    \item Since $D_i$ is independent of $(Y_i(0), Y_i(1))$, $Y_i$ has a linear conditional expectation function:
    $$\mathbb{E}[Y_i \mid D_i] = \mathbb{E}[D_i Y_i(1) + (1 - D_i)Y_i(0) \mid D_i]$$
    $$= D_i\mathbb{E}[Y_i(1)] + (1 - D_i)\mathbb{E}[Y_i(0)]$$
    $$= \mathbb{E}[Y_i(0)] + D_i \times \tau_{ATE}$$
    \item We can run a linear regression of $Y_i$ on $D_i$ and a constant.
    \item The coefficient of $D_i$ is an estimator for $\tau_{ATE}$.
    \item Straightforward calculation shows that this OLS estimator is identical to the difference-in-means estimator.
\end{itemize}
\end{frame}

\begin{frame}{Fisher's Exact Test: The Lady Tasting Tea}
\begin{itemize}
  \item The idea behind Fisher's exact test is often introduced through the famous "lady testing tea" experiment.
  \item Muriel Bristol claimed that she could distinguish cups in which milk was poured first froms cups in which tea was poured first
  \item To test this claim, Fisher prepared 8 cups in random order: 4 milk-first cups and 4 tea-first cups.
  \item After tasting all cups, Bristol was asked to identify the 4 milk-first cups.
  \item The quesiton: How surprsing would her result be if she were simply guessing?
\end{itemize}



\end{frame}
\begin{frame}
\begin{itemize}
    \item Let $\{Y_i(0), Y_i(1)\}_{i=1}^n$ be Bristol's potential answers ($Y_i(d) = 0$ or $1$).
    \item Fisher's randomization scheme is a Completely Randomized Experiment (CRE).
    \item We may interpret
    $$\widehat{\tau}_{DM} = \frac{1}{4} \sum_{i:D_i=1} Y_i - \frac{1}{4} \sum_{i:D_i=0} Y_i = \frac{1}{4}\left(\sum_{i=1}^8 D_i Y_i(1) - \sum_{i=1}^8 (1 - D_i) Y_i(0)\right)$$
    as a measure of her ability.
    \item Because there are only $n = 8$ cups and the outcome is binary, the normal approximation is not appropriate.
\end{itemize}

\end{frame}
\begin{frame}
\begin{itemize}
    \item Consider the sharp null hypothesis
    $$H_0 : Y_i(1) = Y_i(0), \forall i.$$
    \item This is the strongest form of the no treatment effect hypothesis, and differs from the hypothesis that the average treatment effect is zero.
    \item Under $H_0$, $\{Y_i(0), Y_i(1)\}_{i=1}^n$ is known conditional on $\{Y_i\}_{i=1}^n$.
    \item Thus, conditional on $\{Y_i\}_{i=1}^n$, we can compute the exact distribution
    $$\widehat{\tau}_{DM} = \frac{1}{4} \sum_{i=1}^n D_i Y_i(1) - \frac{1}{4} \sum_{i=1}^n (1 - D_i) Y_i(0).$$
\end{itemize}

\end{frame}


\begin{frame}
Under
$$H_0 : \quad \forall i, \quad Y_i(1) = Y_i(0),$$
we have
\pause
\begin{center}
\begin{tabular}{ccccccccc}
\hline
$i$ & 1 & 2 & 3 & 4 & 5 & 6 & 7 & 8 \\
\hline
$Y_i$ & $Y_1$ & $Y_2$ & $Y_3$ & $Y_4$ & $Y_5$ & $Y_6$ & $Y_7$ & $Y_8$ \\
$Y_i(1)$ & $Y_1$ & $Y_2$ & $Y_3$ & $Y_4$ & $Y_5$ & $Y_6$ & $Y_7$ & $Y_8$ \\
$Y_i(0)$ & $Y_1$ & $Y_2$ & $Y_3$ & $Y_4$ & $Y_5$ & $Y_6$ & $Y_7$ & $Y_8$ \\
$D_i$ & 1 & 0 & 1 & 1 & 0 & 0 & 1 & 0 \\
\hline
\end{tabular}
\end{center}

\begin{itemize}
    \item Conditional on $\{Y_i\}_{i=1}^n$, $D$ can take $\binom{8}{4} = 70$ different combinations.
\end{itemize}
\end{frame}

\begin{frame}{Fisher's Exact Test: Test Statistic}

\begin{itemize}
    \item Under the sharp null hypothesis,
    $$H_0 : Y_i(1) = Y_i(0) \quad \text{for all } i,$$
    all potential outcomes are known.
    \item We can therefore recompute the treatment effect estimate for every possible treatment assignment allowed by the randomization scheme.
    \item For each assignment, compute the test statistic
    $$T_n = \frac{\widehat{\tau}_{DM}}{\sqrt{\widehat{V}_{\tau}}}, \quad \text{if } \widehat{V}_{\tau} > 0.$$
    \item The randomization distribution of $T_n$ is the exact distribution induced by the assignment mechanism under $H_0$.
\end{itemize}

\end{frame}

\begin{frame}{Fisher's Exact Test: Exact Validity}

\begin{itemize}
    \item Let $c_{1-\alpha}$ be the $(1-\alpha)$-quantile of the randomization distribution of $|T_n|$ under $H_0$.
    \item Fisher's exact test rejects when
    $$|T_n^{obs}| > c_{1-\alpha}.$$
    \item By construction, the test has conditional level $\alpha$:
    $$\mathbb{P}\left(|T_n| > c_{1-\alpha} \mid \{Y_i\}_{i=1}^n\right) \le \alpha.$$
    \item Conditional validity implies unconditional validity by the law of iterated expectations.
    \item Important: Fisher's exact test requires the sharp null $Y_i(1) = Y_i(0)$ for every unit.
\end{itemize}


\end{frame}

\begin{frame}{Interpretation of Asymptotics}

From now on, we treat the observed sample as drawn from a large population, rather than adopting Fisher's finite-population framework.

Why Does This Matter?
\begin{itemize}
    \item For most statistical methods, this choice does not significantly affect the results.
    \item This perspective aligns more naturally with real-world data applications.
    \item It simplifies reasoning about sampling variability.
\end{itemize}

Both perspectives are valid.
\end{frame}










\begin{frame}{Short Summary}
So far we have already know what is the RCT and how to solve the proof.


Now I want to present you the detailed assignment mechanism.

In short, we knew the "Bernoulli Trial" in the previous slides. Then we are going to show more formal proofs.

Again, I would like to emphasize the assignment mechanism. It means I should follow such mechanism to determine who is treated and who is in control group. 

Example: \\
\textbf{Mechanism}: Imagine walking down a line of 100 people. For person A, you flip a coin. For person B, you flip the coin again. Whether person A got the treatment has absolutely zero impact on the odds for person B.
\end{frame}

\subsection{Stratified Randomized Experiments}

\begin{frame}

\begin{itemize}
    \item The goal of this course is to move beyond RCTs.
    \item The simplest way to move beyond one RCT is to consider two RCTs.
    \item \textbf{Toy Example:} Suppose that we are interested in giving teenagers cash incentives to discourage them from smoking:
    \begin{itemize}
        \item[$\blacksquare$] A random subset of teenagers in Berlin, and
        \item[$\blacksquare$] a random subset of teenagers in Bonn are eligible for the study.
    \end{itemize}
    \item Within each city, we have an RCT.
\end{itemize}

\end{frame}



\begin{frame}{Simpson's Paradox}

Assume that we have the following outcomes:

\begin{center}
\begin{tabular}{lcc}
\hline
Berlin & Non-Smoker & Smoker \\
\hline
Treatment & 152 & 5 \\
Control & 2362 & 122 \\
\hline
\end{tabular}
\quad
\begin{tabular}{lcc}
\hline
Bonn & Non-Smoker & Smoker \\
\hline
Treatment & 581 & 350 \\
Control & 2278 & 1979 \\
\hline
\end{tabular}
\end{center}

$$\widehat{\tau}_{\text{Berlin}} = \frac{5}{152 + 5} - \frac{122}{2362 + 122} \approx -1.7\%, \quad \widehat{\tau}_{\text{Bonn}} = \frac{350}{350 + 581} - \frac{1979}{2278 + 1979} \approx -8.9\%,$$

\begin{center}
\begin{tabular}{lcc}
\hline
Berlin + Bonn & Non-Smoker & Smoker \\
\hline
Treatment & 733 & 355 \\
Control & 4640 & 2101 \\
\hline
\end{tabular}
\pause
\quad
$$\widehat{\tau}_{\text{BB}} = \frac{355}{733 + 355} - \frac{2101}{4640 + 2101} \approx 1.5\%$$
\end{center}

\begin{itemize}
    \item Within each city, the treatment appears to reduce smoking. However, in the aggregate data, the treatment appears to be harmful $\implies$ \textbf{Simpson's paradox}.
\end{itemize}
\end{frame}

\begin{frame}{Stratified Experiments - Setting}

\begin{itemize}
    \item Suppose that covariates $X_i$ take values in a discrete space $X_i \in \mathcal{X}$ s.t. $|\mathcal{X}| = p < \infty$.
    \item Suppose that the treatment assignment is random conditionally on $X_i$ (i.e., we have an RCT in each group defined by a level of $x$):
    $$(Y_i(0), Y_i(1)) \perp\!\!\!\perp D_i \mid X_i = x, \quad \text{for all } x \in \mathcal{X}.$$
    \item Consider an estimator of the ATE by aggregating estimates of the CATE:
    $$\widehat{\tau}_{\text{STRAT}} = \sum_{x \in \mathcal{X}} \frac{n_x}{n} \widehat{\tau}(x), \quad \widehat{\tau}(x) = \frac{1}{n_{x1}} \sum_{\{i : X_i = x, D_i = 1\}} Y_i - \frac{1}{n_{x0}} \sum_{\{i : X_i = x, D_i = 0\}} Y_i$$
    where $n_x = |\{i : X_i = x\}|$ and $n_{xd} = |\{i : X_i = x, D_i = d\}|$.
    \item Alternative interpretation: Apply the difference-in-means estimator after stratifying the sample using the covariates $X_i$;
\end{itemize}

$\implies$ We refer to it as the \emph{stratified estimator}.
\end{frame}


\begin{frame}{Stratified Experiments - Assumptions}

\begin{definitionbox}[Assumption - Superpopulation]
Suppose that $\{X_i, Y_i(0), Y_i(1), D_i\}_{i=1}^n$ are i.i.d. from some distribution $P$ where $X_i$ takes values in a finite cardinality set $\mathcal{X}$ and $\mathbb{E}[Y_i(d)|X_i] < \infty$ for $d \in \{0, 1\}$.
\end{definitionbox}
\begin{definitionbox}[Assumption - Treatment Variation]
There is non-trivial treatment variation for each $x \in \mathcal{X}$, i.e., writing $e(x) = \mathbb{P}[D_i = 1 \mid X_i = x]$, we have $0 < e(x) < 1$ for all $x$.
\end{definitionbox}
\end{frame}

\begin{frame}{Stratified Experiments - Result}
\begin{theorembox}[Theorem Asymptotic Distribution STRAT]
Under our assumptions, we have that
$$\sqrt{n}(\widehat{\tau}_{\text{STRAT}} - \tau_{ATE}) \xrightarrow{d} \mathcal{N}(0, V_{\text{STRAT}}),$$
$$V_{\text{STRAT}} = \text{Var}[\tau_{CATE}(X_i)] + \mathbb{E}\left[\frac{\sigma_{(1)}^2(X_i)}{e(X_i)} + \frac{\sigma_{(0)}^2(X_i)}{1 - e(X_i)}\right],$$
$$\text{with } \sigma_{(d)}^2(X_i) = \text{Var}(Y_i(d)|X_i) \text{ for } d \in \{0, 1\}$$
\end{theorembox}
The asymptotic variance $V_{\text{STRAT}}$ does not depend on $|\mathcal{X}| = p$, the number of groups, or equivalently the number of "parameters" $\tau(x)$ estimated.
\end{frame}



\begin{frame}


\end{frame}



\begin{frame}{Treatment Assignment Level: Clustering}
\begin{itemize}
  \item So far, we have assumed assignment occurs at individual level
  \item Sometimes, random assignment occurs at cluster level {the range is restricted, school, company...}
  \item Individual assignment is too costly
  \item Treatment only makes sense at cluster level but outcome is measured at individual level (i.e. Teacher training experiment)
  \item All three randomization schemes can easily be modified to the cluster level
  \item For Bernoulli trial and CRE, OLS estimator can be used to estimate ATE and cluster-robust standard errors are advisable.

\end{itemize}
\end{frame}

\subsection{Predetermined Covariates}

\begin{frame}{Regression Adjustment in Randomized Trials}
\begin{itemize}
  \item In RCTs, we often observe additional pre-treatment covariates $X_i$, along with treatment assignment $D_i$ and outcomes $Y_i$
  \item So researchers often use regressions instead fo relying only on the difference-in-means estimator.
  \item I use \textbf{interacted regression:} $Y_i \sim \alpha +\tau D_i+X_i^T\beta +D_i\cdot(X_i^T \gamma)$
  \item Then I try to analyze teh interacted regression estimator $\hat{\tau}_{IREG}$, which:
  \begin{itemize}
    \item Allows for transparent modeling of treatment effect heterogeneity.
    \item Is widely considered best practice in modern causal inference.
  \end{itemize}
\end{itemize}





\end{frame}
\begin{frame}{Regression Adjustments in Randomized Trials}

\begin{itemize}
    \item The ATE can be estimated as the difference in predicted outcomes if everyone were treated versus if no one were treated:
    $$\widehat{\tau}_{\text{IR}} = \frac{1}{n} \sum_{i=1}^n \left(\widehat{\alpha} + \widehat{\tau} + X_i^\top (\widehat{\beta} + \widehat{\gamma})\right) - \frac{1}{n} \sum_{i=1}^n \left(\widehat{\alpha} + X_i^\top \widehat{\beta}\right) = \widehat{\tau} + \left(\frac{1}{n} \sum_{i=1}^n X_i\right)^\top \widehat{\gamma}.$$
    \item The interacted regression estimator is algorithmically equivalent to:
    \begin{enumerate}
        \item Running separate OLS regressions for treated and control units:
        $$Y_i \sim \alpha(0) + X_i \cdot \beta(0) \quad \text{for } D_i = 0 \quad \text{and} \quad Y_i \sim \alpha(1) + X_i \cdot \beta(1) \quad \text{for } D_i = 1$$
        \item computing $\widehat{\tau}_{\text{IREG}} = \widehat{\alpha}(1) - \widehat{\alpha}(0) + \overline{X}^\top (\widehat{\beta}(1) - \widehat{\beta}(0))$
    \end{enumerate}
\end{itemize}

\end{frame}

\begin{frame}{Regression Adjustment Under Linearity}
\begin{itemize}
    \item We first impose the assumption that the linear model is correctly specified, and that $X_i \in \mathbb{R}$, e.g.,
    $$Y_i(d) = \alpha(d) + X_i \cdot \beta(d) + \varepsilon_i(d),$$
    with $\mathbb{E}[\varepsilon_i(d) \mid X_i] = 0, \text{Var}(\varepsilon_i(d) \mid X_i) = \sigma^2$.
    \item Treatment is assigned via a Bernoulli trial, and observed outcomes are:
    $$Y_i = Y_i(D_i) = \alpha(0) + D_i(\alpha(1) - \alpha(0)) + X_i \cdot \beta(0) + D_i X_i \cdot (\beta(1) - \beta(0)) + \varepsilon_i,$$
    with $\mathbb{E}[\varepsilon_i \mid X_i, D_i] = 0, \text{Var}(\varepsilon_i \mid X_i, D_i) = \sigma^2$, and $\pi = 0.5$.
    \item For simplicity, we assume that $\text{Var}(X_i) = A$ and $\mathbb{E}[X_i] = 0$
\end{itemize}

\end{frame}
\begin{frame}{Variance of $\hat{\tau}_{DM}$ and $\tau_{IREG}$ under the Linear Model}





\end{frame}


\begin{frame}{What if without Linear Causal Effect}

If the causal effect is not linear, we need to change our way of estimating.

For example, we need to use nonparametric estimation...

\end{frame}


\begin{frame}





\end{frame}
\begin{frame}{Balancing Checks}
\begin{itemize}
  \item Due to randomization of $D_i$, treatment and control groups should, on average, have similar pre-treatment characterstics.
  \item Test whether treatment status is related to pre-treatment covariates, e.g.$$H_0:\beta_1=0 \quad vs. \quad H_1: \beta_1\neq 0 \quad \text{where} \quad D_i=\beta_0+\beta_1 Z_i+u_i$$
  \item Only use covariates measured "before treatment" or "unaffected" by treatment
  \item Balance checks are diagnostic tools, not proof that the randomization was implemented correctly.
  \item Failure to reject does not verify that randomization was valid, and if we test many covariates,some will appear imbalanced purely by chance.
  \item In practice, it is useful to report both tests of individual and joint significance and \textbf{standardized mean differences}, which often more informative than raw $p$-values.

\end{itemize}

\end{frame}{}
\begin{frame}












\end{frame}

\section{Critique to RCT}
\begin{frame}{RCTs Are Powerful, But Not Always Sufficient}
\begin{itemize}
  \item Randomized controlled trials are the benchmark for causal inference.
  \item But they do not solve every question.
  \item Some reasons:
\begin{enumerate}
  \item some treatments cannot be randomized,
  \item experiments may be imperfectly implemented,
  \item results from one experiment may not carry over to other settings ro larger-scale policies.
\end{enumerate}
\item RCTs are extremely valuable, but they are not a complete solution to all empirical questions.
\end{itemize}

\end{frame}
\begin{frame}{Some Treatments Cannot Be Randomized}
\begin{itemize}
  \item Unethical, illegal, or impractical. Example: we cannot randomly assign individuals to smoke, to experience unemployment, or to receive a recession.
  \item As a result, many important causal questions cannot be answered with an RCT, and we have to rely on observational methods.
\end{itemize}




\end{frame}
\begin{frame}{Two Questions}
Once an experiment is conducted, we should ask two separate questions:
\begin{enumerate}
  \item Can we learn the causal effect of the study population?
  \item Can we use this result beyond the study population and setting?

\end{enumerate}

These correspond to:
\begin{itemize}
  \item Internal validity: is the estimated effect credible for the units actually studied?
  \item External validity: does this effect generalize to other populations, settings, or policy environments?
\end{itemize}

\end{frame}
\begin{frame}{Threats to Internal Validity}

True experiments isolate causality by keeping treatment assignment completely independent of potential outcomes.


\begin{itemize}
  \item Flawed Randomization: The actual assignment deviates from the planned design.
  \item Non-adherence: Assigned treatment and received treatment do not match.
  \item Attrition Bias: Data goes missing unevenly across the experimental arms.
\end{itemize}

\end{frame}


\begin{frame}{Imperfect Compliance}
\begin{itemize}
  \item In some experiments, treatment assignment does not coincide with actual treatment.
  \item Example: College Tuition Scholarship
  \begin{itemize}
    \item Treated individuals are randomly offered a substantial scholarship to attend university.
    \item Control individuals are not offered the scholarship.
    \item Some individuals who are offered the scholarship still choose not to enroll in university (e.g., because they decide to enter the workforce immediately).
    \item Some control individuals still enroll in university by paying out-of-pocket or taking on student debt.
  \end{itemize}
\end{itemize}
  \item Thus, assignment and treatment status differ.
  \item Also a topic in instrumental variables.
\end{frame}
\begin{frame}{External Validity}
\begin{itemize}
  \item Even if an experiment is internally valid, its results may not generalize.
  \item External validity asks whether the estimated effect applies to other populations, settings, or policy environments.
  \item Example: education intervention in different countries may have a different effect.
  \item When interpreting experimental evidence, ask:
  \begin{itemize}
  \item Who was studied?
  \item In which setting?
  \item Would the same effect be plausible elsewhere?
\end{itemize}
\end{itemize}


\end{frame}

\begin{frame}{From Small Experiments to Large-Scale Policies}
\begin{itemize}
  \item An experiment often measures the effect of a treatment in small and controlled setting.
  \item This is a partial equilibrium effect: the broader environment is held fixed.
  \item When a policy is scaled up, prices, behavior, and institutions may adjust.
  \item In addition, spillovers may arise: one unit's treatment can effect another unit's outcome.
\end{itemize}

\end{frame}



\section{Partial Identification}

\begin{frame}

I will only introduce a small part of this chapter, because I will then upload a video to discuss partial identification. Note that I will also mention this definition in the following courses.

\end{frame}

\begin{frame}


\end{frame}

\begin{frame}


\end{frame}

\begin{frame}


\end{frame}

\begin{frame}


\end{frame}

\begin{frame}


\end{frame}

\begin{frame}


\end{frame}





\section{Key Takeaways}

\end{document}




